\section{Dirichlet series and factorization}\label{sec:dirichlet-series}

In this section, we introduce the two-variable Dirichlet series
$F(s, y)$ and establish its key analytic properties. The main
result is the factorization
$F(s, y) = \zeta(s)^{2+2y} H(s, y)$, where $H(s, y)$ is holomorphic
and uniformly bounded for $s$ near 1 and $y$ near 0.

\subsection{Euler product}\label{sec:euler-product}

Recall the definition
\begin{equation}
    F(s, y) := \sum_{n \geq 1} \frac{d(n)(1+y)^{\omega(n)}}{n^s}.
\end{equation}
Since the arithmetic function $n \mapsto d(n)(1+y)^{\omega(n)}$ is
multiplicative in $n$ for each fixed $y$, the series admits an Euler
product representation for $\Real(s) > 1$~\cite{Apostol1976}.

\begin{lemma}\label{lem:euler-product-conv}
Fix $\eta\in(0,1/2)$ and choose an integer $K=K(\eta)$ with
$K\geq3$. Then, for every $y$ with $|y|\leq\eta$,
\[
    \bigl|d(n)(1+y)^{\omega(n)}\bigr|\leq d_K(n)
    \qquad(n\geq1),
\]
where $d_K(n)$ is defined by
\[
    \zeta(s)^K=\sum_{n\geq1}d_K(n)n^{-s}
    \qquad (\Real(s)>1).
\]
Consequently, the Dirichlet series $F(s,y)$
converges absolutely for $\Real(s)>1$, uniformly for $|y|\leq\eta$.
\end{lemma}

\begin{proof}
Put $M=1+\eta$. Since $|y|\leq\eta$,
\[
    \bigl|d(n)(1+y)^{\omega(n)}\bigr|
        = d(n)\,|1+y|^{\omega(n)}
        \leq d(n)\,M^{\omega(n)}.
\]
Both sides are multiplicative in $n$. For a prime power $p^a$ with $a\geq1$,
\[
    d(p^a)M^{\omega(p^a)}=(a+1)M,
    \qquad
    d_K(p^a)=\binom{a+K-1}{K-1}.
\]
The quantity $\frac{1}{a+1}\binom{a+K-1}{K-1}$ is increasing in $a$
for $K\geq3$, and at $a=1$ it equals $K/2$. Since $K/2>M$, we obtain $(a+1)M\leq d_K(p^a)$ for every
$a\geq1$. Extending multiplicatively yields
$d(n)M^{\omega(n)}\leq d_K(n)$, and hence the claimed domination.

Since $\zeta(s)^K$ converges absolutely for $\Real(s)>1$, the
domination gives the asserted absolute convergence of $F(s,y)$,
uniformly for $|y|\leq\eta$.
\end{proof}


\begin{lemma}\label{lem:euler-product}
For $\Real(s) > 1$ and every $y \in \C$,
\begin{equation}\label{eq:euler-product}
    F(s, y) = \prod_p \left(1 + (1+y)\left((1 - p^{-s})^{-2} - 1\right)\right),
\end{equation}
where the product is over all primes $p$.
\end{lemma}

\begin{proof}
By multiplicativity, for $\Real(s) > 1$ we have
\[
    F(s, y) =\sum_{n \geq 1}\frac{d(n)(1+y)^{\omega(n)}}{n^s} = \prod_p \left(\sum_{k \geq 0}
    \frac{d(p^k)(1+y)^{\omega(p^k)}}{p^{ks}}\right).
\]
For $k = 0$, the summand equals $1$. For $k \geq 1$, we have
$d(p^k) = k + 1$ and $\omega(p^k) = 1$, so the local factor becomes
\[
    1 + (1+y) \sum_{k \geq 1} \frac{k+1}{p^{ks}}.
\]
Setting $u := p^{-s}$ and using the standard identity
\[
    \sum_{k \geq 1} (k+1) u^k = \frac{1}{(1-u)^2} - 1
    \quad \text{for } |u| < 1,
\]
the local factor equals $1 + (1+y)\bigl((1-p^{-s})^{-2} - 1\bigr)$,
which proves~\eqref{eq:euler-product}.
\end{proof}

\subsection{Factorization isolating the singularity}\label{sec:factorization}

The Riemann zeta function admits the Euler product~\cite{Apostol1976}
$\zeta(s) = \prod_p (1 - p^{-s})^{-1}$ for $\Real(s) > 1$. We define the corresponding power through its local
factors:
\[
    \zeta(s)^{2+2y}
        := \prod_p (1 - p^{-s})^{-2-2y},
    \qquad
    (1 - p^{-s})^{-2-2y}
        := \exp\bigl((-2-2y)\,\log(1 - p^{-s})\bigr),
\]
where $\log$ denotes the principal branch of the logarithm.
This is well-defined because $|p^{-s}| < 1$ for $\Real(s) > 1$, so
$1 - p^{-s}$ lies in the open disk of radius $1$ centered at $1$,
which avoids the branch cut $(-\infty, 0]$ of $\log$.
With this convention, each local factor is analytic in $s$ 
for $\Real(s) > 1$ and in $y$ near $0$.

Define
\begin{equation}\label{eq:H-def}
    H(s, y) := \prod_p (1 - p^{-s})^{2+2y}
    \left(1 + (1+y)\left((1 - p^{-s})^{-2} - 1\right)\right).
\end{equation}

\begin{corollary}\label{thm:factorization}
For $\Real(s) > 1$ and $y$ in a neighborhood of $0$,
\begin{equation}\label{eq:factorization}
    F(s, y) = \zeta(s)^{2+2y} H(s, y).
\end{equation}
\end{corollary}

\begin{proof}
This follows by combining Lemma~\ref{lem:euler-product} with the
Euler product for $\zeta(s)^{2+2y}$ and regrouping factors.
\end{proof}

The purpose of the factorization~\eqref{eq:factorization} is to
isolate the singular behavior of $F(s, y)$ at $s = 1$ in the
factor $\zeta(s)^{2+2y}$, leaving $H(s, y)$ as a well-behaved
function in a larger region, as we now establish.

\subsection{Analytic properties of \texorpdfstring{$H(s, y)$}{H(s,y)}}\label{sec:H-analytic}

We show that the Euler product~\eqref{eq:H-def} defines a
holomorphic function on a strip containing $\Real(s) = 1$,
uniformly for $y$ in a small disk around $0$. The key is that
each local factor $H_p(s, y)$ deviates from $1$ by a quantity of
order $p^{-2\Real(s)}$.

\begin{lemma}\label{lem:local-factor-estimate}
Fix $\eta \in (0, 1/2)$ and $\delta \in (0, 1/2)$. Write
$u := p^{-s}$ and denote the local factor by
\begin{equation}\label{eq:Hp-def}
    H_p(s, y) := (1 - u)^{2+2y}\left(1 + (1+y)\left((1-u)^{-2} - 1\right)\right).
\end{equation}
Then for $|y| \leq \eta$ and $\Real(s) \geq 1 - \delta$,
\begin{equation}\label{eq:local-factor-estimate}
    H_p(s, y) = 1 + O_{\eta, \delta}(|u|^2),
\end{equation}
where the implied constant depends only on $\eta$ and $\delta$.
\end{lemma}

\begin{proof}
The condition $\Real(s) \geq 1 - \delta$ implies that, for every
prime $p \geq 2$,
\[
    |u| = p^{-\Real(s)} \leq p^{-(1 - \delta)} \leq 2^{-(1-\delta)} < 1,
\]
since $\delta < 1/2$. In particular, $|u|$ is bounded away from
$1$, and the power series expansions
\[
    (1 - u)^{-2} - 1 = 2u + 3u^2 + 4u^3 + \cdots = 2u + O_\delta(|u|^2),
\]
\[
    (1 - u)^{2+2y} = 1 - (2+2y)u + \binom{2+2y}{2}u^2 + \cdots
                   = 1 - (2+2y)u + O_{\eta, \delta}(|u|^2),
\]
converge absolutely, with remainders uniform for $|y| \leq \eta$
and $\Real(s) \geq 1 - \delta$. Multiplying the expansions:
\[
    H_p(s, y) = \bigl(1 - (2+2y)u + O_{\eta, \delta}(|u|^2)\bigr)
                \bigl(1 + (1+y)(2u + O_\delta(|u|^2))\bigr).
\]
The coefficient of $u^1$ on the right is $-(2+2y) + 2(1+y) = 0$.
Hence $H_p(s, y) = 1 + O_{\eta, \delta}(|u|^2)$,
proving~\eqref{eq:local-factor-estimate}.
\end{proof}

% =====================================================
% Proposition 3.4 (revised): H is jointly holomorphic
% =====================================================

\begin{proposition}\label{prop:H-analytic}
Fix $\eta \in (0, 1/2)$ and $\delta \in (0, 1/2)$, and let
\begin{equation}\label{eq:H-region}
    \mathcal{R}_{\delta, \eta} := \{(s, y) \in \C^2 :
    \Real(s) \geq 1 - \delta, \ |y| \leq \eta\}.
\end{equation}
Then the partial products of the Euler product~\eqref{eq:H-def}
converge absolutely and uniformly on $\mathcal{R}_{\delta, \eta}$
to a function $H(s, y)$. For every $y$ with $|y| \leq \eta$, the
function $H(s, y)$ is holomorphic in $s$ on the open strip
$\Real(s) > 1 - \delta$. Moreover, $H(s, y)$ is jointly holomorphic
in $(s, y)$ on $\operatorname{int}(\mathcal{R}_{\delta, \eta})$,
and $H$ is uniformly bounded on $\mathcal{R}_{\delta, \eta}$.
\end{proposition}

\begin{proof}
By Lemma~\ref{lem:local-factor-estimate}, for every
$(s, y) \in \mathcal{R}_{\delta, \eta}$ and every prime $p$,
\begin{equation}\label{eq:Hp-bound}
    |H_p(s, y) - 1| \leq C_{\eta, \delta} \, p^{-2(1-\delta)},
\end{equation}
where the constant $C_{\eta, \delta}$ depends only on $\eta$ and
$\delta$. Since $\delta < 1/2$, we have $2(1 - \delta) > 1$, and
the series $\sum_n n^{-2(1-\delta)}$ converges; in particular,
$\sum_p p^{-2(1-\delta)}$ converges. By the Weierstrass $M$-test,
the series
\[
    \sum_p |H_p(s,y)-1|
\]
converges uniformly on $\mathcal{R}_{\delta, \eta}$. Hence the
infinite product criterion applies: if
$\sum_n |u_n|$ converges uniformly, then the product
$\prod_n(1+u_n)$ converges absolutely and uniformly. Therefore the infinite product
$\prod_{p} H_p(s, y)$ converge absolutely and uniformly on
$\mathcal{R}_{\delta, \eta}$. In particular, the partial products $\prod_{p\leq P} H_p(s,y)$ converge uniformly to a limit, which we denote by $H(s, y)$.



To deduce holomorphy in $s$, fix $y_0 \in \C$ with $|y_0| \leq \eta$
and let $s_0 \in \C$ be any point with $\Real(s_0) > 1 - \delta$.
Choose $r > 0$ small enough that $r < \Real(s_0) - (1 - \delta)$,
and define
\begin{equation}\label{eq:K-disk}
    K := \overline{D}(s_0, r)
       = \{s \in \C : |s - s_0| \leq r\}.
\end{equation}
The set $K$ is closed and bounded in $\C$, so by the Heine--Borel
theorem it is compact. By the choice of $r$, every $s \in K$
satisfies
\[
    \Real(s) \geq \Real(s_0) - r > 1 - \delta,
\]
and hence $K \times \{y_0\} \subset \mathcal{R}_{\delta, \eta}$.
The uniform convergence of the partial products on
$\mathcal{R}_{\delta, \eta}$ established above therefore restricts
to uniform convergence of $\prod_{p \leq P} H_p(s, y_0)$ on $K$
as $P \to \infty$.

Each finite partial product $\prod_{p \leq P} H_p(s, y_0)$ is
holomorphic in $s$ on the open disk
$D(s_0, r) = \operatorname{int}(K)$, since each factor
$H_p(\cdot, y_0)$ is. By the theorem on uniform limits of
holomorphic functions (Weierstrass), the limit $H(\cdot, y_0)$ is
holomorphic on $D(s_0, r)$, and in particular at $s_0$. Since
$s_0$ with $\Real(s_0) > 1 - \delta$ was arbitrary, $H(\cdot, y_0)$
is holomorphic on the open strip $\Real(s) > 1 - \delta$. Since
$y_0$ with $|y_0| \leq \eta$ was also arbitrary, the holomorphy
claim in $s$ is proved.

For joint holomorphy, each finite partial product
$\prod_{p \leq P} H_p(s, y)$ is holomorphic in $(s, y)$ on
$\operatorname{int}(\mathcal{R}_{\delta, \eta})$, since each factor
$H_p(s, y)$ is. The uniform convergence of the partial products on
$\mathcal{R}_{\delta, \eta}$ established above restricts to uniform
convergence on every compact subset of
$\operatorname{int}(\mathcal{R}_{\delta, \eta})$, since
$\operatorname{int}(\mathcal{R}_{\delta, \eta}) \subset
\mathcal{R}_{\delta, \eta}$. By the theorem on uniform limits of
holomorphic functions in two complex variables, the limit
$H(s, y)$ is jointly holomorphic on
$\operatorname{int}(\mathcal{R}_{\delta, \eta})$.

Uniform boundedness on $\mathcal{R}_{\delta, \eta}$ follows
from~\eqref{eq:Hp-bound}:
\[
    |H(s, y)|
    \leq \prod_p \bigl(1 + |H_p(s, y) - 1|\bigr)
    \leq \exp\left(\sum_p |H_p(s, y) - 1|\right)
    \leq \exp\left(C_{\eta, \delta} \sum_p p^{-2(1-\delta)}\right),
\]
the right-hand side being finite and independent of
$(s, y) \in \mathcal{R}_{\delta, \eta}$.
\end{proof}


% =====================================================
% Lemma 3.5: H(s, 0) = 1
% =====================================================

\begin{lemma}\label{prop:H-at-zero}
The function $H(s, y)$ satisfies
\begin{equation}\label{eq:H-at-y-zero}
    H(s, 0) \equiv 1 \quad \text{on } \Real(s) > 1/2.
\end{equation}
In particular, $H(1, 0) = 1$.
\end{lemma}

\begin{proof}
On $\Real(s) > 1$, the factorization~\eqref{eq:factorization}
specialized at $y = 0$ reads
\[
    F(s, 0) = \zeta(s)^2 \, H(s, 0).
\]
On the other hand, the Dirichlet series defining $F$ gives
\[
    F(s, 0) = \sum_{n \geq 1} \frac{d(n)}{n^s} = \zeta(s)^2
    \quad \text{on } \Real(s) > 1.
\]
Comparing the two expressions yields
$\zeta(s)^2 H(s, 0) = \zeta(s)^2$ on $\Real(s) > 1$. Since
$\zeta(s) \neq 0$ on $\Real(s) > 1$ (which follows directly from
the Euler product representation of $\zeta$), we may divide both
sides by $\zeta(s)^2$ to obtain
\[
    H(s, 0) \equiv 1 \quad \text{on } \Real(s) > 1.
\]

By Proposition~\ref{prop:H-analytic}, for each $\delta \in (0, 1/2)$
the Euler product defines a holomorphic function on
$\Real(s) > 1 - \delta$. Since these are uniform limits of the same
partial products for different $\delta$, they form a single
holomorphic function $H(s, y)$ on $\Real(s) > 1/2$.
In particular $H(s, 0)$ is holomorphic on the connected open set $\Real(s) > 1/2$.
The constant function $1$ is also holomorphic
there, and the two functions agree on the open subset
$\Real(s) > 1$. By the identity theorem, they agree on all of
$\Real(s) > 1/2$. In particular, $H(1, 0) = 1$.
\end{proof}

